% Page 34 du document source
% Transcription provisoire — à relire et corriger si nécessaire.

\begin{center}
    {\Large\itshape Alban au XV\textsuperscript{ème} siècle (suite)}
\end{center}

\vspace{1.5em}

En 1691, le curé de Notre Dame d'Alban, « Mestre » (Docteur en
théologie) Jean Alary, recteur, sentant sans doute sa fin prochaine,
céda le 13 mars 1691, les revenus de cette chapellenie dont il avait la
jouissance, à Pierre Amig, prêtre habitant Bazuéjouls (ou Bazuéjoulet ?)
à charge naturellement d'assurer deux messes par semaine.

\medskip

Le curé Jean Alary qui, en 1671 (il était arrivé en 1669), avait fait
graver, sans doute au-dessus de son presbytère, la devise : \textsc{A Dieu
seul}, avec un cœur servant de trait d'union à ces deux mots (\textit{Soli
Deo}), mourut effectivement quelque temps après avoir fait cette
donation. La pierre portant gravés son nom et sa devise a été conservée
et se trouve sous le banc de pierre, à droite, devant la façade de
l'église. On peut peut-être signaler que la présence d'un cœur dans la
devise s'expliquerait assez facilement, s'il y avait là un écho des
apparitions du Sacré Cœur à Sainte Marguerite Marie. Effectivement, les
révélations faites à la Sainte sont de ce temps-là.
